\exercisesection{5}

¶ 1. 设 A 为整环，K 为其分式域，T 为 A 上的线性拓扑（第三章 § 2，习题 21）。

\begin{enumerate}
\def\labelenumi{(\alph{enumi})}
\item
  要使 \(\mathcal{T}\) 下 0 的邻域构成某个与 K 的环结构相容的拓扑 \(\mathcal{T}_{\mathrm{K}}\) 在 0 处的基本邻域系，其充要条件是 \(\mathcal{T}\) 为拓扑 \(\mathcal{T}_{u}(\mathrm{A})\)（第三章 §2，习题 24）；此时相对于 \(\mathcal{T}_{\mathrm{K}}\)，A 是有界子集，而 \(\mathcal{T}_{\mathrm{K}}\) 是局部有界的 Hausdorff 拓扑（一般拓扑，第三章 §6，习题 12 与 20 (e)）。K 关于 \(\mathcal{T}_{\mathrm{K}}\) 完备（分别线性紧、严格线性紧（第三章 §2，习题 21））的充要条件是 A 关于 \(\mathcal{T}\) 具有相应性质。
\item
  拓扑 （其中 \(\mathcal{T}_{\mathbf{K}}\)）与 \(\mathcal{T} = \mathcal{T}_u(\mathbf{A})\) 的域结构相容的充要条件是 \(\mathbf{K}\) 的 Jacobson 根  满足 \(\Re(\mathbf{A})\)。\(\mathbf{A}\)\(\neq 0\)。
\end{enumerate}

¶ 2. 设 K 为（交换）域，T 为 K 上与其环结构相容的非离散 Hausdorff 拓扑。拓扑 T 由 K 上的赋值或绝对值定义的充要条件是 T 局部逆有界（一般拓扑，第三章 § 6，习题 22）。若 K 中存在拓扑幂零元，使用一般拓扑第三章 § 6 的习题 22 (d) 以及一般拓扑第九章 § 3 的习题 13。反之则使用一般拓扑第三章 § 6 的习题 22 (f)。

¶ 3. 设 A 为 Noether 整环，K 为其分式域，\(\mathcal{T}_{u}\) 为拓扑 \(\mathcal{T}_{u}(\mathrm{A})\)，\(\mathcal{T}_{\mathrm{K}}\) 为 K 上的相应拓扑（习题 1）。

\begin{enumerate}
\def\labelenumi{(\alph{enumi})}
\item
  若 A 是 Jacobson 根为 \(\mathfrak{r} \neq 0\) 的 Zariski 环，且 A 关于 r-进拓扑完备，证明 A 关于拓扑 \(\mathcal{T}_u\) 完备（一般拓扑，第三章 § 3，第 5 节，命题 9 的推论 2）。
\item
  假设  非离散且由 \(\mathcal{T}_{\mathbf{K}}\) 上的赋值 \(v\)\(\mathbf{K}\) 定义。证明 v 是离散赋值，且 \(v\) 是 v 的环 \(\mathbf{A}\) 的开子环；并证明逆命题。（利用蕴含 \(\mathbf{V}\) 的假设以及 §4 第 1 节的命题 1，可设 \(v\)\(\mathbf{A} \neq \mathbf{K}\)。利用 \(\mathbf{A} \subset \mathbf{V}\) 在 \(\mathbf{A}\) 下是开的事实，证明 \(\mathcal{T}_{\mathbf{K}}\) 是有限生成的 A-模，并使用 §3，第 5 节的命题 9 得出结论。对逆命题，注意 \(\mathbf{V}\) 是局部环，其中极大理想 \(\mathbf{A}\) 与 (0) 是唯一的素理想，因而其中每个理想 \(m\) 都是 \(a \neq 0\)-准素理想。此时 \(m\) 的整闭包为 \(\mathbf{A}\)。给出 \(\mathbf{V}\) 的例子（令 \(\mathbf{A} \neq \mathbf{V}\) 为形式幂级数环 \(\mathbf{V}\)，其中 k 为域）。\(k[[T]]\)\(k\)，其中 k 为域）。
\item
  推出一个非离散完备拓扑域的例子，其拓扑不是局部逆有界的（使用 (b) 和习题 2）。
\end{enumerate}

\begin{enumerate}
\def\labelenumi{\arabic{enumi}.}
\setcounter{enumi}{3}
\tightlist
\item
  设  为域 K 上的赋值，A 为 v 的环，m 为其极大理想。\(v\)\(K\)\(v\)\(m\) 为域 K 上的赋值，A 为 v 的环，m 为其极大理想。
\end{enumerate}

\begin{enumerate}
\def\labelenumi{(\alph{enumi})}
\item
  A 上由 v 定义的拓扑与 m-进拓扑相同的充要条件是 A 为域或离散赋值环。
\item
  由  定义的拓扑使 A 成为严格线性紧环（第三章 §2，习题 21）的充要条件是 A 为域或离散赋值环（使用 (a) 及第三章 §2，习题 22 (a)）。\(v\)\(A\)\(A\) 定义的拓扑使 A 成为严格线性紧环（第三章 §2，习题 21）的充要条件是 A 为域或离散赋值环（使用 (a) 及第三章 §2，习题 22 (a)）。
\end{enumerate}

\begin{enumerate}
\def\labelenumi{\arabic{enumi}.}
\setcounter{enumi}{4}
\tightlist
\item
  \begin{enumerate}
  \def\labelenumii{(\alph{enumii})}
  \tightlist
  \item
    设 K 为域，v 为 K 上的赋值，\(\Gamma\) 为 v 的阶群。K 关于拓扑 \(T_{v}\) 线性紧（第三章 §2，习题 15）的充要条件是：对 \(\Gamma\) 的每个良序子集 B 以及 K 的每个元素族 \((a_{\beta})_{\beta\in\mathcal{B}}\)，若对 \(\lambda<\mu<v\)，
  \end{enumerate}
\end{enumerate}

\[
v (a _ {\lambda} - a _ {\mu}) <   v (a _ {\mu} - a _ {\nu}),
\]

则存在元素 ，使得对满足 \(a \in \mathbf{K}\) 的每个指标有序对 ，都有 。（使用集合论第三章 § 2 的习题 4，\(v(a - a_{\lambda}) = v(a_{\lambda} - a_{\mu})\)\(\lambda, \mu\)\(\mu > \lambda\)。（使用集合论第三章 § 2 的习题 4，

第三章 § 2）。若此条件成立，则赋值 v 的环 A 关于离散拓扑也线性紧。\(v\)第三章 § 2）。若此条件成立，则赋值 v 的环 A 关于离散拓扑也线性紧。

\begin{enumerate}
\def\labelenumi{(\alph{enumi})}
\setcounter{enumi}{1}
\tightlist
\item
  证明 §3，习题 2 中定义的域  关于 \(S(\Gamma, k) = K\) 线性紧。令 \(T_{v}\) 为有理数的全序群 Q，并考虑 K 的子环 \(\Gamma\)，它由满足下述条件的 \(K_{0}\) 组成：使 \(x = (x_{\alpha})\) 的  所成集合是有限集，或是趋于 \(\alpha \in Q\) 的严格递增序列的点集。证明 \(x_{\alpha} \neq 0\)\(+\infty\) 是域，关于 K 上的典范赋值 v 所诱导的赋值 \(K_{0}\) 完备，并且 \(v_{0}\) 与 v 具有相同的取值群和剩余域，但 \(v_{0}\) 不是线性紧的（参见 §10，习题 2）。\(K_{0}\) 不是线性紧的（参见 §10，习题 2）。
\end{enumerate}

¶ 6. (a) 设 A 为赋值环，M 为 A-模，M' 为 M 的子模。证明 M' 为 M 的纯子模（第一章 § 2，习题 24）的充要条件是：对所有 \(\alpha \in A\)，M' ∩ (\(\alpha\) M) = \(\alpha\) M'。

\begin{enumerate}
\def\labelenumi{(\alph{enumi})}
\setcounter{enumi}{1}
\item
  设  是 \(\mathbf{M}'\) 的纯子模，并且若在 \(\mathbf{M}\) 上赋予这样的拓扑：\(\mathbf{M}'\)（其中 \(\alpha\mathbf{M}'\)）构成 0 的基本邻域系，则 \(\alpha \in \mathbf{A}, \alpha \neq 0\) 线性紧。证明对每个元素 \(0\)\(\mathbf{M}'\)，存在 \(x \in \mathbf{M}\)，使得 \(y' \in \mathbf{M}'\) 具有如下性质：对所有满足 \(x' = x + y'\) 的 ，都有 \(\alpha \in \mathbf{A}\)。（对于每个满足  的 \(x + \mathbf{M}' \in \alpha(\mathbf{M} / \mathbf{M}')\)，考虑 \(x' \in \alpha\mathbf{M}\)\(\alpha \in \mathbf{A}\) 的子集 ，它由满足 \(x + \mathbf{M}' \in \alpha(\mathbf{M} / \mathbf{M}')\)\(\mathrm{S}_{\alpha}\) 的  组成。）\(\mathbf{M}'\)\(y_{\alpha}'\)\(x + y_{\alpha}' \in \alpha\mathbf{M}\) 组成。）
\item
  设 A 为线性紧赋值环；令 K 为 A 的分式域。证明：若 M 是无挠 A-模，且 \(\mathbf{M}_{(\mathbf{K})}\) 有可数基，则 M 是一族可数个秩 1 的 A-模的直和。（将 M 看作递增的纯子模序列 \((\mathbf{M}_{i}^{\prime})\) 的并，其中 \(M_{i}^{\prime}\) 的秩为 i，并对每个 i 将 (b) 应用于 \(M_{i}^{\prime}\) 及其子模 \(M_{i-1}^{\prime}\)。）
\item
  在 (c) 中关于 A 与 M 的假设下，令 N 为 M 的有限秩纯子模。证明 N 是 M 的直因子（注意，每个作为有限个秩 1 的 A-模之直和的 A-模都是线性紧的，并使用 (b)）。
\item
  在 (b) 中关于 M 与  的相同假设下，证明对所有 \(M'\)，存在 \(x \in M\)，使得 \(y' \in M'\) 的湮没子等于元素 \(x' = x + y'\) 的湮没子。（令 a 为 \(x + M' \in M/M'\) 的湮没子；对所有 \(x + M'\)，令 \(\alpha \in a\) 为 M' 的子集，它由满足 \(T_{\alpha}\) 的  组成；证明 \(M'\)\(y'\) 的交非空。）\(\alpha(x + y') = 0\)\(T_{\alpha}\) 的交非空。）
\item
  设 A 为赋值环，并满足对每个理想 ，A-模 A/a（带离散拓扑）线性紧。证明每个有限生成的无挠 A-模 M 都是有限个单生成 A-模的直和。（若 \(a \neq 0\) 是 M 的一组生成元，则对 n 作归纳，取一个湮没子最小的 \((z_{i})_{1 \leqslant i \leqslant n}\)，并注意它生成的 M 的子模是纯的。）\(z_{i}\)，并注意它生成的 M 的子模是纯的。）
\end{enumerate}

\begin{enumerate}
\def\labelenumi{\arabic{enumi}.}
\setcounter{enumi}{6}
\tightlist
\item
  设 K 为域，v 为 K 上的离散赋值，且 K 关于 \(T_{v}\) 完备；令 A 与 p 分别表示 v 的环与理想，U = A - p 表示 A 的可逆元所成的集合。令 u 为 K 到 K 的一个子域上的同构。
\end{enumerate}

\begin{enumerate}
\def\labelenumi{(\alph{enumi})}
\item
  证明  且 \(u(\mathfrak{p}) \subset \mathbf{A}\)。（为证明第一点，注意对所有 \(u(\mathbf{U}) \subset \mathbf{U}\)，当 n > 0 与 v 的剩余域的特征互素时，方程 \(z \in \mathfrak{p}\) 由 Hensel 引理在 \(x^n = 1 + z\) 中有解；若 \(\mathbf{A}\)，则推出整数 \(n > 0\)\(v\)\(v(u(z)) < 0\) 将被每个与 v 的剩余域的特征互素的整数 \(v(u(z))\) 整除。然后由第一点推出第二个断言。）\(n > 0\)\(v\) 整除。然后由第一点推出第二个断言。）
\item
  从 (a) 推出：或者 ，或者 u 连续（考虑 v 的一个统一化元在 u 下的像）。\(u(\mathbf{K}^*) \subset \mathbf{U}\)\(u\)\(u\)\(v\)，或者 u 连续（考虑 v 的一个统一化元在 u 下的像）。
\item
  给出一个代数闭域 \(\Omega\) 的例子，使得取 \(\mathbf{K} = \Omega((\mathbf{X}))\) 时，\(\mathbf{K}\) 同构于 \(\mathbf{K}\) 的一个包含于 \(\mathbf{U} \cup \{0\}\) 的子域（参见代数，第 V 章 § 5，习题 13）。
\end{enumerate}

¶ 8. (a) 设 K 为域，v 为 K 上的高度 1 赋值，A 为其环。令 H 为 K 的紧子集（带拓扑 \(T_{v}\)），令 \(a \neq 0\) 为 K 的一点。证明存在无常数项的多项式 \(f \in K[X]\)，使得 \(f(a) = 1\)，\(f(H) \subset A\)。（证明 f 可取为如下形式的多项式

\[
1 - (1 - a ^ {- 1} \mathrm{X}) (1 - c _ {1} ^ {- 1} \mathrm{X}) ^ {n (1)} \dots (1 - c _ {r} ^ {- 1} \mathrm{X}) ^ {n (r)}
\]

其中  是  中适当选取的满足 \(c_{i}\)\(\mathbf{H}\) 的元素，而  是充分大的整数 \(v(c_{i}) < v(a)\)\(n(i)\)。 \(>0\)。 

\begin{enumerate}
\def\labelenumi{(\alph{enumi})}
\setcounter{enumi}{1}
\tightlist
\item
  令 X 为全不连通紧空间；赋予从 X 到 K 的连续映射环 \(\mathcal{C}(\mathbf{X};\mathbf{K})\) 以一致收敛拓扑。令 B 为 \(\mathcal{C}(\mathbf{X};\mathbf{K})\) 的包含常数且能分离 X 中各点的子环；证明 B 在 \(\mathcal{C}(\mathbf{X};\mathbf{K})\) 中稠密（使用 (a) 及一般拓扑，第 X 章 §4，习题 21 (b)）。
\end{enumerate}

¶ 9. 设 A 为离散赋值环，π 为 A 的统一化元，K 为 A 的分式域；假设 A 关于 π-进拓扑完备。内射 A-模（代数，第二章 §2，习题 11）与可除 A-模（代数，第七章 §2，习题 3）相同，并且都是同构于 K 或 K/A 的 A-模的直和（代数，第七章 §2，习题 3）。此外，每个单生成 A-模都同构于 K/A 的一个子模。A-模 Hom \(_{A}\) (M, K/A) 称为 A-模 M 的（代数）托里对偶，记作 M*；典范映射 c \(_{M}\) : M → M** 是单射（代数，第二章 §2，习题 13 (b)）。对于 M 的每个子模 N，M* 的由满足 u(x) = 0 的 u 组成的子模 N \(^{0}\) 称为 N 在 M* 中的正交补；M/N 的对偶典范地识别为 N \(^{0}\)，N 的对偶识别为 M*/N \(^{0}\)。直极限 lim M \(_{\alpha}\) 的托里对偶典范同构于逆极限 lim M \(_{\alpha}^{*}\)。

\begin{enumerate}
\def\labelenumi{(\alph{enumi})}
\tightlist
\item
  我们知道，秩 1 的 A-模（代数，第七章 § 4，
\end{enumerate}

习题 22）同构于以下形式之一的模：A、K、K/A 或 A/ \(\pi^{h}\) A。证明 K 与 A/ \(\pi^{h}\) A 的代数托里对偶分别同构于 K 与 A/ \(\pi^{h}\) A，A 的托里对偶同构于 K/A，而 K/A 的托里对偶同构于 A（利用 K 的 A-子模的知识以及 A 完备这一事实（相对于 K/A 的对偶））。推出：对每个有限秩 A-模 M，M* 是具有相同秩的模，且典范同态 \(c_{M}\) 为双射。

\begin{enumerate}
\def\labelenumi{(\alph{enumi})}
\setcounter{enumi}{1}
\item
  设 M 为 A-模，N 为 M 的有限秩子模。证明  中  的正交补通过 \(N^{0}\) 识别为 N（使用 (a)）。\(M^{**}\)\(c_{M}\) 识别为 N（使用 (a)）。
\item
  证明 Noether（分别 Artin）A-模 M 具有有限秩（将 M 嵌入其内射包中（代数，第二章 §2，习题 18））；于是 M* 是 Artin（分别 Noether）模。
\item
  令  表示 M* 上的拓扑，即 M 上的逐点收敛拓扑（赋予 K/A 离散拓扑）。证明：若 N 是 M 的子模，则 \(\sigma(M^{*}, M)\) 是由 \(M^{*}\)\(M\)\(K/A\)\(N\)\(M\)\(\sigma(N^{0}, M/N)\) 在  上诱导的拓扑，而 \(N^{0}\) 是拓扑  关于 \(\sigma(M^{*}, M)\) 的商拓扑。（关于第二点，注意若 P 是 M 的满足 \(\sigma(M^{*}/N^{0}, N)\)\(N^{0}\) 的子模，则 P/N 的对偶识别为 \(\sigma(M^{*}, M)\)\(P\)\(M\)\(N \subset P\)。）若 \(P/N\)\(N^{0}/P^{0}\)，则拓扑 \(M = \lim M_{\alpha}\) 是各拓扑 \(\sigma(M^{*}, M)\) 的逆极限。\(\sigma(M_{\alpha}^{*}, \overrightarrow{M_{\alpha}})\) 的逆极限。
\item
  拓扑  与 \(\sigma(\mathbf{K},\mathbf{K})\) 是 \(\sigma(\mathbf{A},\mathbf{K}/\mathbf{A})\)-进拓扑；拓扑 \(\pi\) 与 \(\sigma(\mathbf{A}/\pi^{h}\mathbf{A},\mathbf{A}/\pi^{h}\mathbf{A})\) 是离散拓扑。推出对每个 A-模 M，带有 \(\sigma(\mathbf{K}/\mathbf{A},\mathbf{A})\) 的模  是线性紧的（第三章 §2，习题 15；将 M 看作自由 A-模的商模）。\(\mathbf{M}^{*}\)\(\sigma(\mathbf{M}^{*},\mathbf{M})\) 是线性紧的（第三章 §2，习题 15；将 M 看作自由 A-模的商模）。
\item
  令 M、N 为两个 A-模；对于每个同态 ，令 \(u: \mathbf{M} \to \mathbf{N}\) 表示从 \(^{t}u\) 到 \(\operatorname{Hom}(u, \mathbf{l}_{\mathbf{K/A}})\)\(\mathbf{N}^{*}\) 的同态 \(\mathbf{M}^{*}\)，满足，满足
\end{enumerate}

\[
\left(^ {t} u (w)\right) (x) = w (u (x))
\]

对所有  与所有 \(x \in \mathbf{M}\)，证明 \(w \in \mathbf{N}^*\) 关于拓扑 \({}^t u\) 与 \(\sigma(\mathbf{N}^*, \mathbf{N})\) 连续。若 u 是 \(\sigma(\mathbf{M}^*, \mathbf{M})\) 的自同态 ，则 \(u\)\(x \mapsto \pi x\) 是 \(\mathbf{M}\)\({}^t u\) 的自同态 。对 \(w \mapsto \pi w\) 的每个子模 ，都有 \(\mathbf{M}^*\)\(\mathbf{P}\)。\(\mathbf{M}\)\((u(\mathbf{P}))^0 = {}^t u^{-1}(\mathbf{P}^0)\)。

¶ 10. 假设与记号同习题 9。

\begin{enumerate}
\def\labelenumi{(\alph{enumi})}
\item
  若 M 是离散拓扑 A-模，证明 M 是无挠模。若进一步 M 线性紧，证明 M 是 Artin 模（若 N 是自同态 \(x \mapsto \pi x\) 的核，注意 N 线性紧且离散，因而可看作向量 (A/ \(\pi\) A)-空间；然后使用第二章 § 2 的习题 20 (d)）。
\item
  由 (a) 推出每个线性紧拓扑 A-模都是严格线性紧的（第二章 § 2，习题 19）。
\item
  拓扑 A-模 M 的拓扑托里对偶是代数托里对偶  的子模 ，由从 M 到 K/A 的连续同态组成（后者赋予离散拓扑）。若 M 上的拓扑是线性的（第二章 § 2，习题 14），且 N 是 M 的闭子模，证明 M/N 的拓扑托里对偶识别为 \(M'\)，而 N 的拓扑托里对偶识别为 \(M^{*}\)\(N^{0} \cap M'\)（为确定 N 的对偶，注意对从 N 到 K/A 的每个连续同态 u，存在 M 的开子模 U，使得 u(x) = 0 对 \(M' / (N^{0} \cap M')\) 中的 x 成立，然后使用习题 9）。\(u(x) = 0\)\(N \cap U\) 中的 x 成立，然后使用习题 9）。
\item
  对于有限秩拓扑 A-模 M，拓扑托里对偶 \(M'\) 等于代数托里对偶 \(M^{*}\)（归结为秩 1 的模的情形）。
\item
  设 M 为拓扑线性的 Hausdorff 拓扑 A-模；证明对 M 中所有 ，存在 \(x \neq 0\) 使得 \(u \in M'\)（注意存在 M 的开子模 U 使得 \(u(x) \neq 0\)）；换言之，典范映射 \(x \notin U\) 是单射。推出若 N 是 M 的闭子模，则 \(\mathbf{M} \to (\mathbf{M}')^*\) 关于拓扑 \(N^0 \cap M'\) 在  中稠密（使用 (d)），从而在 \(N^0\) 中有 。\(\sigma(\mathbf{M}^*, \mathbf{M})\)\(N = M \cap (N^0 \cap M')^0\)\(M^{**}\)。
\end{enumerate}

证明 M 关于拓扑  在  中稠密。\((\mathbf{M}^{\prime})^{*}\)\(\sigma((\mathbf{M}^{\prime})^{*}, \mathbf{M}^{\prime})\) 中稠密。

\begin{enumerate}
\def\labelenumi{(\alph{enumi})}
\setcounter{enumi}{5}
\item
  设 M 为线性紧拓扑 A-模；证明典范嵌入 \(M \rightarrow (M')^{*}\) 是双射，且 M 上的拓扑（将 M 识别为 \((M')^{*}\) 后）等于 \(\sigma(M, M')\)。（注意若 U 是 M 的开子模，则由 (a) M/U 是 Artin 模，因而 \(U^{0}\) 是 Noether 模（习题 9 (c)），并推出 M 上所考虑的拓扑相同；借助 (e) 完成证明。）
\item
  设 M、N 为拓扑均为线性的两个拓扑 A-模；对于每个连续同态 ，有 \(u \colon \mathbf{M} \to \mathbf{N}\)，并且也（滥用记号地）用 \(^t u(\mathbf{N}') \subset \mathbf{M}'\) 表示从 \(^t u\) 到 \(\mathbf{N}'\) 的线性映射，其图与 \(\mathbf{M}'\) 相同。若 M 与 N 是 Hausdorff 的，证明 \(^t u\) 在 M 上的限制与 u 重合（分别将 M、N 看作典范嵌入 \(^t (^t u)\) 与 \(u\)\((\mathbf{M}')^*\) 中）；此外，对 N 的每个子模 Q，\((\mathbf{N}')\) 中）；此外，对 N 的每个子模 Q，
\end{enumerate}

\[
(\bar {u} ^ {1} (\mathbf {Q})) ^ {0} \cap \mathbf {M} ^ {\prime} = ^ {t} u (\mathbf {Q} ^ {0})
\]

（使用 (e)）。

\begin{enumerate}
\def\labelenumi{(\alph{enumi})}
\setcounter{enumi}{7}
\tightlist
\item
  设 M 为线性紧 A-模；由 (g) 和习题 9 (f) 推出 M 无挠的充要条件是 M' 可除，而 M 可除的充要条件是 M' 无挠。
\end{enumerate}

